\documentclass[11pt, a4paper]{article}
\usepackage[utf8]{inputenc}
\usepackage[brazil]{babel}
\usepackage[top=2cm, bottom=2cm, left=2cm, right=2cm]{geometry} % Margens ajustadas para caber em 2 páginas
\usepackage{graphicx}
\usepackage{booktabs}
\usepackage{hyperref}
\usepackage{titlesec}
\usepackage{float}
\usepackage{amsmath} % Garante suporte a comandos matemáticos avançados

% Ajuste de espaçamento dos títulos para economizar espaço
\titlespacing*{\section}{0pt}{10pt}{5pt}
\titlespacing*{\subsection}{0pt}{8pt}{4pt}

\title{\vspace{-2cm} \textbf{Relatório do Trabalho Prático 1: Análise de Reviews Steam}}
\author{
    \textbf{Disciplina:} DCC012 - Estrutura de Dados II \\
    \textbf{Professora:} Barbara Quintela
}
\date{Novembro de 2025}

\begin{document}

\maketitle
\thispagestyle{empty} % Remove número da página na primeira página se necessário

\section{Identificação e Detalhamento das Atividades}

Abaixo listam-se os integrantes do grupo e suas respectivas responsabilidades na implementação do trabalho.

\begin{table}[h]
    \centering
    \begin{tabular}{@{}ll@{}}
        \toprule
        \textbf{Integrante (Matrícula)} & \textbf{Atividades Realizadas} \\ \midrule
        Leandro Fortunato (202476031) & Menu, QuickSort, Otimização Arquivos, União. \\ \midrule
        Davi Grossi S. de Souza (202165023A) & Funções \texttt{getReview(int i)} e \texttt{import(int n)}. \\ \midrule
        Rafael Amaral Rodrigues (202165123AB) & Structs Hash, \texttt{createTable} e Sondagem Linear. \\ \midrule
        Yan Dias Ferreira (202335037) & Funções Hash (Divisão/Mult), Benchmarks e Exportação CSV. \\ \midrule
        Rafael Rocha Ribeiro (202365131A) & TAD \texttt{GameReview}, Leitura CSV e Conversão Binário. \\ \bottomrule
    \end{tabular}
\end{table}

\noindent \textbf{Repositório do Projeto:} O código fonte completo pode ser acessado em: \\
\url{https://github.com/Rafaelrr5/TrabED2}

\section{Decisões de Implementação}

O sistema foi desenvolvido na linguagem \textbf{C} para garantir controle eficiente de memória e desempenho, conforme os requisitos de manipulação de arquivos binários. As principais decisões técnicas foram:

\subsection{Implementação das Funções \texttt{getReview} e \texttt{import}}

Como parte da Etapa 1 (Processamento dos Dados) foram implementadas as funções \texttt{getReview(int i)} e \texttt{GameReview* import(int n)}, responsáveis pelo acesso direto ao arquivo binário de reviews e pela importação de reviews aleatórios.

A implementação foi realizada utilizando exclusivamente funções da biblioteca padrão da linguagem C, seguindo o comportamento definido no padrão ISO/IEC 9899:2011 (C11). A referência utilizada encontra-se em: \texttt{https://web.cs.dal.ca/\~vlado/pl/C\_Standard\_2011-n1570.pdf}. Foram empregadas as seguintes funções da biblioteca \texttt{stdio.h}:

\begin{itemize}
    \item \textbf{\texttt{fopen}} (Seção 7.21.5.3) — usada para abrir o arquivo \texttt{reviews.bin} em modo binário;
    \item \textbf{\texttt{fseek}} (Seção 7.21.9.2) — utilizada para posicionar o cursor diretamente no registro desejado, permitindo acesso aleatório em $O(1)$ utilizando \texttt{i * sizeof(GameReview)};
    \item \textbf{\texttt{fread}} (Seção 7.21.8.1) — responsável por ler os bytes correspondentes a um ou mais registros para dentro de estruturas alocadas em memória;
    \item \textbf{\texttt{fprintf}} (Seção 7.21.6.1) — usada apenas em mensagens de depuração durante o desenvolvimento.
\end{itemize}

\subsection{Armazenamento e Acesso aos Dados}
Para cumprir o requisito de acesso aleatório eficiente, os dados do arquivo CSV são convertidos para um arquivo binário (\texttt{reviews.bin}) contendo estruturas de tamanho fixo (\texttt{GameReview}). Isso possibilita o uso da função \texttt{fseek} para acessar o $i$-ésimo registro em tempo constante $O(1)$, evitando o carregamento total do dataset na memória RAM durante a importação de amostras.

\subsection{Estrutura da Tabela Hash}
Utilizou-se uma Tabela Hash com \textbf{endereçamento aberto} e estratégia de \textbf{sondagem linear} para o tratamento de colisões.
\begin{itemize}
    \item \textbf{Justificativa:} A sondagem linear elimina o \textit{overhead} de alocação dinâmica para nós de listas encadeadas (necessário no encadeamento exterior) e favorece a localidade de referência (cache), tornando as iterações mais rápidas.
    \item \textbf{Dimensionamento:} A tabela é inicializada com tamanho $M = 2N + 1$, onde $N$ é o número de registros a serem inseridos. Isso garante um fator de carga $\alpha < 0.5$, condição essencial para manter o desempenho da sondagem linear próximo de $O(1)$.
\end{itemize}

\subsection{Funções de Hashing}
Foram implementadas duas funções para comparação de desempenho:
\begin{enumerate}
    \item \textbf{Método da Divisão (MOD):} $h(k) = |k| \bmod M$. Escolhido pela simplicidade e rapidez de cálculo.
    \item \textbf{Método da Multiplicação:} $h(k) = \lfloor M (k A - \lfloor k A \rfloor) \rfloor$, com $A \approx 0.618$ (razão áurea). Escolhido por sua capacidade de dispersar chaves que apresentam padrões sequenciais, comuns em IDs de banco de dados.
\end{enumerate}

\subsection{Ordenação dos Resultados}
Para exibir os jogos mais avaliados ("Top X"), foi usado um \textbf{QuickSort} ($O(N \log N)$), utilizando a função \texttt{qsort} da biblioteca padrão. O critério de ordenação é a quantidade de reviews (decrescente), utilizando a nota média como desempate.

\section{Análise Comparativa das Funções Hash}

Os testes foram realizados utilizando a opção de benchmark do programa, medindo o tempo e colisões para um conjunto de $N = 1000$ registros.

\begin{table}[H]
    \centering
    \caption{Comparativo de Desempenho (N = 1000 registros)}
    \begin{tabular}{@{}lccc@{}}
        \toprule
        \textbf{Função Hash} & \textbf{Tempo (s)} & \textbf{Colisões} & \textbf{Colisões/Reg} \\ \midrule
        Método da Divisão (MOD)      & 0.018660 & 415 & 0.415 \\
        Método da Multiplicação (MULT) & 0.001075 & 437 & 0.437 \\ \bottomrule
    \end{tabular}
    \label{tab:comparativo}
\end{table}

\textbf{Análise dos Resultados:}
Observou-se que o Método da \textbf{Divisão (MOD)} apresentou um número ligeiramente inferior de colisões (415 contra 437) para a amostra testada. No entanto, o Método da \textbf{Multiplicação} demonstrou um tempo de execução significativamente menor (0.001s contra 0.018s). Embora a diferença absoluta de colisões seja pequena (ambos próximos de 0.4 colisões por registro), a eficiência temporal do método da multiplicação sugere que ele pode ser mais vantajoso em cenários onde a velocidade de cálculo do hash é crítica, mesmo com uma taxa de colisão levemente superior.

\section{Referências Bibliográficas}

\renewcommand{\refname}{} % Remove o título automático "Referências" para economizar espaço
\vspace{-1cm}
\begin{thebibliography}{9}

\bibitem{Cormen2012}
CORMEN, Thomas H. et al. \textit{Algoritmos: teoria e prática}. Tradução de: Vandenberg D. de Souza. 3. ed. Rio de Janeiro: Elsevier, 2012.

\bibitem{Kaggle2025}
CRAINBRAMP. \textit{Steam Dataset 2025: Multi-Modal Gaming Analytics}. Disponível em: https://www.kaggle.com/datasets/crainbramp/steam-dataset-2025-multi-modal-gaming-analytics. Acesso em: nov. 2025.

\end{thebibliography}

\end{document}